\section{Related Work}
\label{sec:related_work}
 
\paragraph{Benchmarks in Minecraft} 
MineDojo~\citep{minedojo} introduces a suite of 1,560 creative tasks defined by natural language instructions. However, it suffers from significant redundancy and overly complex tasks, complicating practical evaluation, as shown in~\cref{fig:mcu_vs_minedojo}. BEDD~\citep{milani2023bedd} defines five tasks covering various aspects of Minecraft, primarily designed for the MineRL BASALT competition. By decomposing the evaluation framework, BEDD facilitates detailed assessments of agent performance across subgoals and attributes such as human likeness. However, its reliance on human ratings limits scalability. Other works on Minecraft agents~\citep{deps,jarvis1,groot,plan4mc,vpt,steve1} lack a unified benchmark, with each agent evaluated on its own task set. We argue that establishing a standardized benchmark is a high priority for advancing Minecraft agent development.

\paragraph{Open-ended Agents in Minecraft} Many agents have been developed to interact with Minecraft environments.
Some methods focus on using Imitation Learning or Reinforcement Learning to learn various skills in open-world Minecraft to complete \textbf{short-horizon tasks}~\citep{vpt,groot,steve1,jiang2025reinforcement,zhao2024optimizing,groot15,yuan2024pre,cai2024groot,jiang2025visual}.
\citet{vpt} generates action labels from Minecraft videos on YouTube using IDM and performs imitation learning to obtain an unconditional policy capable of completing various tasks. \citet{steve1} uses MineCLIP as a text encoder to obtain a text-conditioned multitask policy based on unconditional VPT~\citep{vpt}.
Minecraft also supports the testing of \textbf{long-horizon programmatic tasks}, so some methods accomplish long-horizon tasks by using large language models as planners, combined with skill policies~\citep{deps,minedreamer,chen2024apt,plan4mc,mp5,zheng2023steve}.
\citet{deps} has designed an agent pipeline based on GPT, which interactively completes tasks from environmental feedback through the self-explanation and zero-shot planning capabilities of LLMs.
In order to further enhance the long-horizon capabilities of LLM Agents, some methods have explored efficient explicit memory mechanisms to support agents retrieve and improve from previous interaction trajectories~\citep{jarvis1,gitm,voyager,park2024mr,li2024optimus}.
Unlike designing complex agent pipelines, the rise of VLA~\citep{Brohan2023RT2VM,openvla} has inspired researchers to use end-to-end VLM as a policy to directly fulfill human instructions in Minecraft~\citep{omnijarvis,steve-2}.
Some methods use a code-as-policy approach~\citep{codeaspolicies}, using MineFlayer~\citep{mineflayer} as a language-conditioned policy, combined with a pretrained LLM to accomplish various tasks to avoid agents trapped on short-horizon skills lacking when executing long-horizon tasks~\citep{voyager,octupus,liu2024rl,liu2024odyssey}.
Finally, there are some methods focused on completing \textbf{open-ended creative tasks} in Minecraft, such as building and decoration, which differ from the traditional programmatic object-centric tasks and often cannot be directly defined by rule-based rewards~\citep{guo2024luban,zhang2023creative}.


\paragraph{LLM-as-a-Judge} 
The advancement of large language models (LLMs) has demonstrated remarkable performance in instruction following, query understanding, and response generation. This capability has motivated the use of LLMs as judges~\citep{llm-judge-mtbench}, leveraging their ability to score and rank model outputs. The strong performance of LLMs~\citep{brown2020language}, combined with well-designed assessment pipelines~\citep{li2023prd,beigi2024model,bai2024benchmarking}, enables fine-grained and detailed judgments for various evaluation applications, addressing the limitations of traditional evaluation methods that require extensive human annotation. Recent efforts have also explored the use of LLMs and vision-language models (VLMs) to evaluate AI agents~\citep{mobile-pan,zhuge2024agent}. Compared to previous work, MCU is the first to apply this paradigm to open-ended game agent task generation and evaluation. We argue that LLM-based task generators have the potential to create diverse tasks that are crucial for evaluating open-ended game agents, given the inherent open-endedness of LLMs~\citep{hughes2024open}. Furthermore, using the same LLM (or VLM) as a judge ensures more consistent evaluation criteria compared to crowdsourcing approaches~\citep{web-env-autoagent}.